
\documentclass{article}
\usepackage{colortbl}
\usepackage{makecell}
\usepackage{multirow}
\usepackage{supertabular}

\begin{document}

\newcounter{utterance}

\centering \large Interaction Transcript for game `cladder', experiment `full\_v1.5\_default', episode 3456 with refineonly{-}Qwen3.5{-}4B.
\vspace{24pt}

{ \footnotesize  \setcounter{utterance}{1}
\setlength{\tabcolsep}{0pt}
\begin{supertabular}{c@{$\;$}|p{.15\linewidth}@{}p{.15\linewidth}p{.15\linewidth}p{.15\linewidth}p{.15\linewidth}p{.15\linewidth}}
    \# & \multicolumn{2}{c}{Player} && \multicolumn{2}{c}{Game Master} \\
    \hline

    \theutterance \stepcounter{utterance}  
    & & & \multicolumn{4}{p{0.6\linewidth}}{
        \cellcolor[rgb]{0.9,0.9,0.9}{
            \makecell[{{p{\linewidth}}}]{
                \texttt{\tiny{[P1$\langle$GM]}}
                \texttt{You are an expert in causal inference. The following question is not a typical commonsense query, but rather a meticulously designed question created by a professor specializing in causal inference, intended to assess the students' mastery of the course content.} \\
\texttt{Imagine a self{-}contained, hypothetical world with only the following conditions, and without any unmentioned factors or causal relationships: Solar eclipse has a direct effect on Alice arriving to school. Traffic has a direct effect on Alice arriving to school. Method 1: We look at how solar eclipse correlates with Alice arriving to school case by case according to traffic. Method 2: We look directly at how solar eclipse correlates with Alice arriving to school in general. To understand how solar eclipse affects Alice arriving to school, is it more correct to use the Method 1 than Method 2?} \\
\texttt{Start your answer with 'Yes' or 'No', followed by additional reasoning or evidence to support your explanation.} \\
            }
        }
    }
     \\ \\

    \theutterance \stepcounter{utterance}  
    & \multicolumn{4}{p{0.6\linewidth}}{
        \cellcolor[rgb]{0.561,0.737,0.561}{
            \makecell[{{p{\linewidth}}}]{
                \texttt{\tiny{[P1$\rangle$GM]}}
                \texttt{Yes} \\
\\ 
\texttt{In this hypothetical world, **Traffic** acts as a **confounding variable** (or a common cause) because it has a direct causal effect on Alice arriving to school, while the Solar Eclipse is the variable of interest.} \\
\\ 
\texttt{Here is the step{-}by{-}step causal reasoning:} \\
\\ 
\texttt{1.  **Identify the Structure**: The causal structure is \$Solar Eclipse \textbackslash{}rightarrow Alice\$ and \$Traffic \textbackslash{}rightarrow Alice\$. There is no arrow between Solar Eclipse and Traffic; they are independent causes.} \\
\texttt{2.  **Analyze Method 2 (General Correlation)**: If we simply look at the correlation between Solar Eclipse and Alice without conditioning on Traffic, we are observing the **total effect** of the Solar Eclipse. However, because Traffic is a confounder, the observed correlation in Method 2 will be a mixture of the true causal effect of the Solar Eclipse and any spurious associations introduced by the distribution of Traffic. If Traffic is correlated with the Solar Eclipse (e.g., eclipses happen more often when traffic is heavy due to some external factor not mentioned but implied by the "case by case" analysis), Method 2 will yield a biased estimate. Even if Traffic is independent of the Solar Eclipse, Method 2 measures the total effect, which is technically correct for the *total* impact, but often in causal inference contexts involving confounders, we seek the **direct effect** or a causal estimate that isolates the specific pathway.} \\
\texttt{3.  **Analyze Method 1 (Conditioning on Traffic)**: By looking at the correlation "case by case according to traffic," we are effectively **conditioning on the confounder** (Traffic). This is the definition of **stratification** or **backdoor adjustment**.} \\
\texttt{    *   Mathematically, this calculates \$P(Alice | Solar Eclipse, Traffic)\$.} \\
\texttt{    *   By holding Traffic constant, we block the backdoor path from Solar Eclipse to Alice that goes through Traffic.} \\
\texttt{    *   This isolates the **direct causal effect** of the Solar Eclipse on Alice, removing the influence of the Traffic variable.} \\
\texttt{4.  **Conclusion**: In causal inference, when a variable (Traffic) has a direct effect on the outcome (Alice) and we want to understand the specific effect of another variable (Solar Eclipse), we must control for the confounder to avoid bias. Method 1 performs this control, whereas Method 2 leaves the confounding influence unaddressed. Therefore, to accurately determine how the Solar Eclipse *causes* Alice to arrive (isolating the specific mechanism), Method 1 is the statistically and causally correct approach.} \\
            }
        }
    }
    & & \\ \\

    \theutterance \stepcounter{utterance}  
    & & & \multicolumn{2}{p{0.3\linewidth}}{
        \cellcolor[rgb]{0.9,0.9,0.9}{
            \makecell[{{p{\linewidth}}}]{
                \texttt{\tiny{[GM$|$GM]}}
                \texttt{Yes} \\
            }
        }
    }
    & & \\ \\

    \theutterance \stepcounter{utterance}  
    & & & \multicolumn{2}{p{0.3\linewidth}}{
        \cellcolor[rgb]{0.9,0.9,0.9}{
            \makecell[{{p{\linewidth}}}]{
                \texttt{\tiny{[GM$|$GM]}}
                \texttt{no} \\
            }
        }
    }
    & & \\ \\

    \theutterance \stepcounter{utterance}  
    & & & \multicolumn{2}{p{0.3\linewidth}}{
        \cellcolor[rgb]{0.9,0.9,0.9}{
            \makecell[{{p{\linewidth}}}]{
                \texttt{\tiny{[GM$|$GM]}}
                \texttt{game\_result = LOSE} \\
            }
        }
    }
    & & \\ \\

\end{supertabular}
}

\end{document}
